% This is a conceptual outline following NeurIPS 2025 format
% You would compile this in LaTeX with the official neurips_2025.sty

\documentclass{article}
\usepackage[preprint, nonatbib]{neurips_2025}
\usepackage{amsmath,amssymb,amsfonts}
\usepackage{booktabs}
\usepackage{graphicx}

\usepackage[numbers]{natbib}


\title{Beyond Uniform Computation:\\
Spatially-Adaptive Sparse Diffusion Transformers for Efficient Text-to-Image Synthesis} 

\title{Midterm Project Proposal: Semantically-Aware Diffusion Transfomers for High-Speed, Low-Compute Text-to-Image Synthesis}

\author{
Brody L. Jordan \\
Texas A\&M University \\
431005357 \\ 
\texttt{blj@tamu.edu}
}

\begin{document}

\maketitle

\begin{abstract}
    Diffusion Transformers (DiTs) have emerged as the dominant architecture for
    high-fidelity text-to-image synthesis, yet they suffer from prohibitive computational
    costs due to their uniform token processing paradigm. This survey examines recent
    advances in efficient DiT architectures, categorizing approaches into architectural
    optimizations, token reduction strategies, and hybrid representations. Through critical
    analysis of state-of-the-art methods including SANA, HART, and Dynamic Token Merging
    techniques, we identify a fundamental limitation: current approaches fail to exploit
    the spatial redundancy inherent in visual data. We propose Spatially-Adaptive Sparse
    Diffusion Transformers (SA-SparseDiT), a novel framework that dynamically allocates
    compute based on prompt-driven saliency, enabling 3--5$\times$ speedups while
    maintaining generation quality.
\end{abstract}

\section{Introduction and Motivation}

Text-to-image synthesis has undergone a revolutionary transformation with the advent of
Diffusion Transformers (DiTs). Models such as Stable Diffusion 3 \cite{esser2024scaling}
and FLUX \cite{labs2024flux} demonstrate unprecedented visual quality by replacing
traditional U-Net backbones with scalable transformer architectures. However, this
quality comes at a steep computational cost: DiTs process \textit{every} latent token
with equal intensity, regardless of semantic importance.

The quadratic complexity $O(N^2)$ of self-attention in DiTs becomes the primary
bottleneck when scaling to high resolutions. For a $1024 \times 1024$ image with
$16\times$ VAE compression, the sequence length $N = 4096$, resulting in $\sim$16M
attention operations per layer per timestep. With typical models requiring 50
denoising steps and 20+ transformer layers, total operations exceed $16 \times 10^9$
per image generation—prohibitive for real-time applications.

This survey focuses on \textbf{Spatially-Adaptive Computation} in DiTs: methods that
reduce this burden by selectively processing tokens based on their importance.
We examine 8 high-impact papers from CVPR, NeurIPS, and ICLR (2024--2025) that
attack this problem from three distinct angles: (1) architectural efficiency,
(2) token reduction strategies, and (3) hybrid representations. Our analysis
reveals that while recent works achieve impressive speedups, they fail to
dynamically align computational allocation with the prompt's semantic requirements—a
critical gap we address in our proposed research direction.

\section{Categorization and Taxonomy}

We organize the surveyed literature into three complementary paradigms:

\subsection{Category A: Architectural Efficiency Improvements}
These methods modify the DiT backbone to reduce computational complexity
\textit{without} reducing token count, through attention approximations or
structural changes.

\textbf{Papers:}
\begin{itemize}
    \item \textbf{SANA 1.5} \cite{xie2025sana}: Linear Diffusion Transformers
          with depth-wise convolution attention, reducing complexity from $O(N^2)$ to
          $O(N)$ while maintaining global receptive fields.
    \item \textbf{EDiT} \cite{becker2025edit}: Efficient Diffusion Transformers
          with Linear Compressed Attention using low-rank factorization of attention
          matrices.
\end{itemize}

\subsection{Category B: Token Reduction and Merging}
These approaches explicitly reduce the sequence length $N$ by merging or
pruning redundant tokens during inference.

\textbf{Papers:}
\begin{itemize}
    \item \textbf{Token Merging (ToMe)} \cite{bolya2023token}: Bipartite
          matching-based merging for Vision Transformers, adapted for diffusion models.
    \item \textbf{ToMA} \cite{lu2025toma}: Token Merge with Attention,
          introducing saliency-aware merging strategies specifically for DiTs.
    \item \textbf{Dynamic Chunking DiT} \cite{haridas2026dynamic}: Recent work
          introducing patch-level scheduling based on local information density.
\end{itemize}

\subsection{Category C: Hybrid Representations and Alternative Paradigms}
These works challenge the pure-diffusion paradigm by combining discrete
autoregressive modeling with continuous refinement.

\textbf{Papers:}
\begin{itemize}
    \item \textbf{HART} \cite{tang2024hart}: Hybrid Autoregressive Transformer
          combining discrete tokens for structure with continuous residuals for
          detail, achieving single-pass generation.
    \item \textbf{DDiT} \cite{kim2026ddit}: Dynamic Patch Scheduling using
          learnable region importance scores.
    \item \textbf{Shiva-DiT} \cite{zhang2026shiva}: Residual-Based
          Differentiable Top-$k$ Selection for token pruning.
\end{itemize}

\begin{table}[h]
    \centering
    \caption{Taxonomy of Efficient DiT Methods}
    \label{tab:taxonomy}
    \begin{tabular}{@{}lccc@{}}
        \toprule
        \textbf{Method}  & \textbf{Category}  & \textbf{Speedup} & \textbf{Granularity}       \\
        \midrule
        SANA 1.5         & A: Architectural   & 2--3$\times$     & Global                     \\
        EDiT             & A: Architectural   & 1.5--2$\times$   & Global                     \\
        ToMe             & B: Token Reduction & 2$\times$        & Spatial (static)           \\
        ToMA             & B: Token Reduction & 2.5$\times$      & Spatial (attention-guided) \\
        Dynamic Chunking & B: Token Reduction & 3$\times$        & Patch-level                \\
        Shiva-DiT        & B: Token Reduction & 2--4$\times$     & Token-level (residual)     \\
        HART             & C: Hybrid          & 10$\times$       & Representation-level       \\
        \bottomrule
    \end{tabular}
\end{table}

\section{Deep Dive and Methodological Comparison}

\subsection{Architectural Efficiency: The SANA Approach}

SANA 1.5 \cite{xie2025sana} replaces standard self-attention with
\textbf{Linear Diffusion Transformer} blocks using depth-wise convolutions.
The attention mechanism becomes:

$$\text{Attention}(Q, K, V) = \text{Softmax}\left(\frac{QK^T}{\sqrt{d}}\right)V
    \approx \text{DWConv}(Q) \cdot V$$

While this achieves linear complexity $O(N)$, it \textit{uniformly} reduces
computation across all tokens. Regions requiring fine detail (faces, text)
receive the same ``lightweight'' attention as simple backgrounds, potentially
compromising high-frequency generation.

\textbf{Trade-off:} SANA trades model expressiveness for speed. The linear
approximation cannot capture long-range dependencies as effectively as full
attention, particularly for complex compositional scenes.

\subsection{Token Reduction Strategies: From Static to Dynamic}

\paragraph{ToMe (Static Merging).}
ToMe \cite{bolya2023token} defines ``similarity'' via cosine distance in feature
space, merging the $r$ most similar token pairs at each layer:

$$
    \mathcal{M} = \{(i, j) \mid \text{argsort-top-}r(\cos(x_i, x_j))\}
$$

\textbf{Limitation:} Feature similarity does not correlate with semantic
importance. Two distinct texture patches may have high feature similarity
but require independent processing for photorealism.

\paragraph{ToMA (Attention-Guided Merging).}
ToMA \cite{lu2025toma} improves upon ToMe by using cross-attention maps
$A \in \mathbb{R}^{N \times M}$ (where $M$ is text sequence length):

$$
    s_i = \sum_{j=1}^{M} A_{i,j}
$$

Tokens with low $s_i$ (poor alignment with text) are prioritized for merging.

\textbf{Advancement:} Text-to-vision alignment improves merging decisions.
However, ToMA applies merging \textit{statically} at fixed layers, failing
to adapt to the evolving importance of tokens across denoising timesteps.

\paragraph{DDiT and Dynamic Chunking.}
Recent work \cite{kim2026ddit, haridas2026dynamic} introduces temporal
adaptivity: token importance is recomputed at each denoising step using
learned ``Decider'' modules or patch-level statistics.

\textbf{Gap:} These methods schedule compute at the \textit{patch} or
\textit{chunk} level, not individual tokens. Coarse-grained decisions
miss opportunities for fine-grained efficiency.

\paragraph{Shiva-DiT.}
Shiva-DiT \cite{zhang2026shiva} implements differentiable Top-$k$ selection
using the Gumbel-Softmax trick:

$$
    P(\text{keep}_i) = \text{Softmax}\left(\frac{\log \pi_i + g_i}{\tau}\right)
$$

where $\pi_i$ is a learned importance score and $g_i \sim \text{Gumbel}(0,1)$.

\textbf{Innovation:} End-to-end trainable sparsity. However, the importance
score $\pi_i$ is learned from reconstruction loss alone, lacking explicit
\textit{text-prompt guidance}.

\subsection{Hybrid Representations: HART}

HART \cite{tang2024hart} abandons the diffusion paradigm entirely for a
hybrid autoregressive approach:

\begin{enumerate}
    \item Discrete VQ-VAE tokens capture global structure in one AR pass.
    \item Continuous residuals add detail via small diffusion heads.
\end{enumerate}

\textbf{Strengths:} Achieves single-pass generation (massive speedup),
leveraging the efficiency of AR models (GPT-style next-token prediction).

\textbf{Fundamental Limitation:} HART (and all AR models) require
\textit{discrete tokenization}, which introduces an information bottleneck
at the VQ-VAE stage. Complex textures and continuous gradients suffer from
quantization artifacts. Furthermore, AR models autoregressively decode,
making them inherently sequential—batch generation of multiple images
cannot be parallelized as effectively as diffusion.

\section{Limitations and Open Challenges}

Our survey reveals three systematic failures across the state-of-the-art:

\subsection{Challenge 1: The Uniformity Assumption}
\textbf{Current approaches assume uniform importance across the spatial
    domain.} Whether applying full attention (standard DiT), linear attention
(SANA), or token merging (ToMe), existing methods process the ``sky''
region with the same intensity as a ``human face.''

\textbf{Why it fails:} Visual scenes exhibit extreme spatial redundancy.
In most images, 70--80\% of pixels are low-entropy (smooth regions,
repetitive textures). Current architectures waste compute on these
``uninformative'' regions while under-investing in high-frequency details
critical for perceived quality.

\subsection{Challenge 2: Static vs. Dynamic Importance}
\textbf{Text-prompt alignment is underutilized in compute allocation.}
While ToMA uses cross-attention for merging decisions, it treats this
saliency as \textit{constant} across timesteps.

\textbf{Why it fails:} The importance of a token evolves during
denoising. At $t=T$ (pure noise), all tokens are equally uncertain.
At $t=T/2$, some regions (e.g., backgrounds) may stabilize while others
(e.g., foreground objects) remain noisy. Current methods lack mechanisms
to dynamically re-prioritize tokens as the diffusion trajectory unfolds.

\subsection{Challenge 3: The Differentiable Trilemma}
\textbf{End-to-end trainable sparsity requires balancing three conflicting
    goals:} (1) differentiability (for gradient-based training), (2) efficiency
(for wall-clock speedup), and (3) quality preservation (no generation
artifacts).

Shiva-DiT achieves (1) and (2) but compromises on (3)—learned importance
scores drift toward conservative selection, reducing effective sparsity.
Feature caching methods achieve (2) and (3) but are \textit{not}
differentiable, requiring hand-tuned heuristics.

\textbf{Why it matters:} Without a unified solution to this trilemma,
spatial sparsity remains a ``post-hoc'' optimization rather than a
first-class architectural primitive.

\subsection{Challenge 4: Seamless Integration}
\textbf{No existing method handles the boundary between sparse and dense
    regions.} When merging tokens or dropping computation, current approaches
risk visible artifacts at region boundaries—the ``blocky'' effect where
high-detail subjects meet low-detail backgrounds.

\textbf{Real-world impact:} These limitations prevent deployment of
high-quality DiTs on consumer hardware. Current models require 24GB+ VRAM
and 10--30 seconds per image, limiting applications in real-time creative
tools, mobile AR, and autonomous systems requiring rapid visual
imagination.

\section{Proposed Research Direction (Towards a NeurIPS Submission)}

\subsection{The Specific Limitation to Solve}
We target \textbf{Challenge 2 (Dynamic Importance)} and \textbf{Challenge 4
    (Seamless Integration)}. Our hypothesis is that \textit{compute allocation
    should be a function of both the text prompt (semantic saliency) and the
    current timestep (temporal stability).}

\subsection{Proposed Solution: SA-SparseDiT}

We propose \textbf{Spatially-Adaptive Sparse Diffusion Transformers
    (SA-SparseDiT)}, a DiT variant with three novel components:

\subsubsection{1. Tri-Signal Saliency Router}
A differentiable gating module that computes per-token priority $P_i$ using
three signals:

$$P_i = \underbrace{\alpha \cdot A_{\text{text}, i}}_{\text{Text Alignment}} +
    \underbrace{\beta \cdot \|\nabla_\theta \epsilon_i\|_2}_{\text{Noise Magnitude}} +
    \underbrace{\gamma \cdot \mathcal{V}(x_i)}_{\text{Spatial Variance}}$$

where:
\begin{itemize}
    \item $A_{\text{text}, i}$: Cross-attention weight to text embeddings
          (semantic importance).
    \item $\|\nabla_\epsilon\|_2$: Magnitude of predicted noise (temporal
          uncertainty).
    \item $\mathcal{V}(x_i)$: Local spatial variance (structural complexity).
\end{itemize}

\textbf{Why this is novel:} Unlike ToMA (static attention) or Shiva-DiT
(residual only), our router dynamically recomputes priorities at each
denoising step, allowing tokens to ``graduate'' from sparse to dense
processing as they become semantically relevant.

\subsubsection{2. Differentiable Top-$k$ Selection via Gumbel-Softmax}
Using the concrete distribution \cite{jang2016categorical}, we implement
sparse selection:

$$
    z_i = \text{sigmoid}\left(\frac{\log P_i + g_i}{\tau}\right), \quad g_i \sim \text{Gumbel}(0,1)
$$

\textbf{Top-$k$ tokens} ($z_i > \tau_{\text{sparsity}}$) proceed through
full self-attention. \textbf{Bottom-$(N-k)$ tokens} are merged into $m \ll k$
anchor tokens via a learned clustering operation.

\subsubsection{3. Seamless Re-Integration via Residual Flows}
To address Challenge 4, we introduce \textbf{seamless un-pooling} using
normalizing flows \cite{dinh2016density}. When expanding merged tokens
back to full resolution, we apply invertible transformations that
preserve local smoothness while allowing high-frequency detail injection
from dense regions.

\subsection{Expected Contributions}
1. \textbf{First prompt-driven dynamic sparsity} for DiTs with per-token,
per-timestep re-prioritization.
2. \textbf{Differentiable sparse attention kernels} enabling end-to-end
training without quality degradation.
3. **3--5$\times$ speedup** on 1024p generation with $<5\%$ FID increase,
enabling 8K synthesis on consumer GPUs.

\subsection{Risk Mitigation}
\begin{itemize}
    \item \textbf{Rebuttal: ``HART already solves efficiency.''}
          $\rightarrow$ HART requires discrete tokenization. SA-SparseDiT
          preserves the continuous flow-matching paradigm, maintaining superior
          texture quality.
    \item \textbf{Rebuttal: ``Dynamic Chunking DDiT already token-prunes.''}
          $\rightarrow$ DDiT operates at patch-level ($16\times16$ blocks).
          SA-SparseDiT is token-level, offering finer granularity.
\end{itemize}

\begin{thebibliography}{99}

    \bibitem{esser2024scaling}
    Patrick Esser et al.
    \newblock Scaling Rectified Flow Transformers for High-Resolution Image Synthesis.
    \newblock \emph{arXiv preprint arXiv:2403.03206}, 2024.

    \bibitem{labs2024flux}
    Black Forest Labs.
    \newblock FLUX.1: A Series of State-of-the-Art Diffusion Models.
    \newblock \emph{Technical Report}, 2024.

    \bibitem{xie2025sana}
    Enze Xie et al.
    \newblock SANA 1.5: Efficient Scaling of Training-Time and Inference-Time Compute in Linear Diffusion Transformer.
    \newblock \emph{arXiv preprint arXiv:2501.18427}, 2025.

    \bibitem{becker2025edit}
    Philipp Becker et al.
    \newblock EDiT: Efficient Diffusion Transformers with Linear Compressed Attention.
    \newblock \emph{arXiv preprint arXiv:2503.16726}, 2025.

    \bibitem{bolya2023token}
    Daniel Bolya and Judy Hoffman.
    \newblock Token Merging for Fast Stable Diffusion.
    \newblock \emph{arXiv preprint arXiv:2303.17604}, 2023.

    \bibitem{lu2025toma}
    Wenbo Lu et al.
    \newblock ToMA: Token Merge with Attention for Diffusion Models.
    \newblock \emph{arXiv preprint arXiv:2509.10918}, 2025.

    \bibitem{haridas2026dynamic}
    Akash Haridas et al.
    \newblock Dynamic Chunking Diffusion Transformer.
    \newblock \emph{arXiv preprint arXiv:2603.06351}, 2026.

    \bibitem{tang2024hart}
    Haotian Tang et al.
    \newblock HART: Efficient Visual Generation with Hybrid Autoregressive Transformer.
    \newblock \emph{arXiv preprint arXiv:2410.10812}, 2024.

    \bibitem{kim2026ddit}
    Dahye Kim et al.
    \newblock DDiT: Dynamic Patch Scheduling for Efficient Diffusion Transformers.
    \newblock \emph{arXiv preprint arXiv:2602.16968}, 2026.

    \bibitem{zhang2026shiva}
    Jiaji Zhang et al.
    \newblock Shiva-DiT: Residual-Based Differentiable Top-k Selection for Efficient Diffusion Transformers.
    \newblock \emph{arXiv preprint arXiv:2602.05605}, 2026.

    \bibitem{jang2016categorical}
    Eric Jang, Shixiang Gu, and Ben Poole.
    \newblock Categorical Reparameterization with Gumbel-Softmax.
    \newblock \emph{ICLR}, 2017.

    \bibitem{dinh2016density}
    Laurent Dinh, Jascha Sohl-Dickstein, and Samy Bengio.
    \newblock Density estimation using Real NVP.
    \newblock \emph{ICLR}, 2017.

\end{thebibliography}

\end{document}
